\scanpage{238}\mspage{\unsure{93}}
Donc pour tt \uline{co}faisceau \inserted{$\unsure{E|\mathcal E}$} sur $V$, à valeurs dans
une catégorie $M$ avec $\underset{\unsure{-|\longrightarrow}}{\lim}$, on a iso
\[
\underset{\unsure{U|u}}{\varinjlim}\,g^*(E)\stackrel{\sim}{\longrightarrow}
\underset{V}{\underset{\unsure{-|\longrightarrow}}{\lim}}\,E,
\]
et pour tt préfaisceau \inserted{$F$} sur $V$, à val. dans une
catégorie $M$ avec $\varprojlim$, on a iso
\[
\underset{U}{\varprojlim}\,g^*(F)\stackrel{\sim}{\longleftarrow}
\underset{V}{\varprojlim}\,F.
\]

Démonstration. \deleted{\unsure{Comme}} \inserted{Je vais \deleted{\unsure{ici}} d'abord expliciter}
\deleted{les morphismes étales sont loc. cofiltrants}

\uline{Lemme} \inserted{7.5.} Soit $h:X'\unsure{-|\to}X$ un morphisme étale
dans $\operatorname{Cat}$. Alors $h$ est loc. cofiltrant. \deleted{\unsure{car en}}\inserted{De façon}
plus précise, pour tout $x$ dans $X$, \deleted{$\,{}_{x\backslash}X'$ est
\unsure{composé} \inserted{\unsure{somme}} \unsure{directe} de ses composantes connexes,}
on a
\[
\pi_0({}_{x\backslash}X')\overset{\sim}{\longrightarrow}\operatorname{Ob} X'_x
\]
toute composante connexe de $\,{}_{x\backslash}X'$ contient
un élément et un seul de \deleted{$\unsure{X'|X'_x}$} la
cat. discrète $X'_x\subset{}_{x\backslash}X'$. Si $x'$ est dans
$X'_x$ (i.e. \deleted{\unsure{avec}} $x'$ objet de $X'$ tel que $h(x')=x$),
c'est un objet \uline{initial} de la composante
connexe qui le contient. (Donc cette composante est cofiltrante.)

Ceci \unsure{prouve}\deleted{\unsure{donc}} \deleted{\unsure{la (ad}} (\unsure{admis cela|admet cela}) la première
assertion de 7.4., compte tenu de 7.3.
(stabilité par composition), on \unsure{prouve|perçoit} qu'une
% archive sheet number: 238
